\documentclass[11pt]{article}

\usepackage[T1]{fontenc}
\usepackage[utf8]{inputenc}
\usepackage{lmodern}
\usepackage{amsmath,amssymb}
\usepackage[margin=0.85in]{geometry}
\usepackage{booktabs}
\usepackage{longtable}
\usepackage{tabularx}
\usepackage{enumitem}
\usepackage{xcolor}
\usepackage[hidelinks]{hyperref}
\usepackage{xurl}

\setlist{leftmargin=*,itemsep=0.25em,topsep=0.35em}
\setlength{\parindent}{0pt}
\setlength{\parskip}{0.55em}
\setlength{\emergencystretch}{3em}

\definecolor{donegreen}{RGB}{30,115,65}
\definecolor{partialamber}{RGB}{160,95,0}
\definecolor{blockred}{RGB}{160,35,35}
\definecolor{externalblue}{RGB}{25,75,145}

\newcommand{\done}{\textcolor{donegreen}{\textbf{Implemented}}}
\newcommand{\partialstatus}{\textcolor{partialamber}{\textbf{Partial}}}
\newcommand{\blocked}{\textcolor{blockred}{\textbf{Submission blocker}}}
\newcommand{\external}{\textcolor{externalblue}{\textbf{External run required}}}
\newcommand{\optional}{\textcolor{externalblue}{\textbf{Optional}}}
\newcommand{\file}[1]{\path{#1}}

\begin{document}

\begin{center}
  {\Large\bfseries Submission-Readiness and Regeneration Plan}\par
  \vspace{0.35em}
  {\large Hybrid Quantum-Classical ADMM--QAOA Manuscript}\par
  \vspace{0.35em}
  Prepared July 24, 2026
\end{center}

\section*{Purpose and recommended submission path}

This document records the remaining work needed to make the corrections to
review findings 1--5 appear consistently in the code, evidence artifacts,
figures, tables, numerical prose, and final submission package.  It is an
execution and acceptance plan, not a source of numerical results.

The recommended path is a \emph{conservative, software-and-simulation
submission}: regenerate the predeclared post-contingency benchmark,
sensitivity analysis, ADMM histories, fixed-topology SOC diagnostics, and
nonlinear AC recovery in the documented environment, but keep quantitative
hardware claims excluded unless a repeated multi-date hardware campaign is
completed.  This path is scientifically complete enough to answer the five
review findings without making a two-job hardware comparison carry more
weight than it can support.

No new post-contingency number should be copied into \file{main_2.tex} by
hand.  Numerical prose must be generated from validated artifacts or be
limited to qualitative statements that remain true independently of a run.

\section*{1. Current corrected baseline}

\begin{longtable}{p{0.13\textwidth}p{0.24\textwidth}p{0.54\textwidth}}
\toprule
\textbf{Review item} & \textbf{Current status} & \textbf{What is already in place} \\
\midrule
\endfirsthead
\toprule
\textbf{Review item} & \textbf{Current status} & \textbf{What is already in place} \\
\midrule
\endhead
1: QUBO structure & \done & Component-level QUBO accounting, the exact
sorting/prefix baseline for the uniform-cardinality special case, exact
radial enumeration, and a flow-based spanning-tree MILP are implemented.
The manuscript no longer equates density with hardness. \\
\addlinespace
2: artifact contradictions & \done{} for the retained artifact; \partialstatus{}
for new results & One strict legacy JSON artifact drives its CSV, figures,
macros, and appendix table.  The old ``success probability'' is correctly
described as an optimizer-run exact-QUBO hit fraction.  Unsupported SA,
ADMM, library-baseline, and hardware values are not reconstructed. \\
\addlinespace
3: continuous model & \done{} in code and equations; \external{} for numbers
& Continuous-model-v2 includes two-sided voltage-drop deactivation,
fixed-power-factor shedding, branch-specific ratings, explicit transformer
taps, slack balance, units, and validation diagnostics.  The retained
continuous-model-v1 losses are marked stale. \\
\addlinespace
4: AC validation & \done{} in validation logic; \external{} for evidence &
The code distinguishes SOC feasibility, cone tightness, angle recovery, and
nonlinear fixed-injection AC power-flow validation.  The manuscript no longer
calls a SOC relaxation result ``AC-true.'' \\
\addlinespace
5: radiality mathematics & \done{} in formulation and audit; \external{} for
the prospective comparison & The historical pairwise cycle and
anti-islanding terms are explicitly called heuristic surrogates.  Exact
spanning-tree constraints, exhaustive false-positive/false-negative auditing,
and cycle-safe final projection with an action log are available. \\
\bottomrule
\end{longtable}

The current retained benchmark remains a base-case, legacy scaling artifact.
It is useful for audit history, but it cannot support a post-contingency
engineering conclusion or a current-model physical comparison.

\section*{2. Non-negotiable evidence hierarchy}

For the final paper, evidence should flow in one direction:
\[
\begin{aligned}
\text{protocol and code}
&\longrightarrow \text{strict JSON artifacts}
\longrightarrow \text{validated manifest}\\
&\longrightarrow \text{generated tables, figures, and macros}
\longrightarrow \text{manuscript prose}.
\end{aligned}
\]

The following rules should be enforced.

\begin{enumerate}
  \item \file{koshi_admm_qaoa/results/study_protocol.json} is frozen before
  prospective execution.  Its hash must be stored in every prospective
  artifact.
  \item Raw or per-seed outputs are never replaced by aggregate-only rows.
  Every stochastic result must preserve its seed, bitstring, objective,
  topology flags, projection record, and continuous-validation record.
  \item JSON must be strict: no \texttt{NaN}, \texttt{Infinity}, or implicit
  missing-value substitutions.  Missing evidence is represented explicitly
  with a documented status and reason.
  \item Generated TeX and PNG files are read-only products of the pipeline.
  If a number must change, the JSON-producing run or generator must change.
  \item The manifest must hash the protocol, primary results, sensitivity
  results, ADMM histories, generated tables, generated macros, and figures.
  \item A clean checkout must be able to run all validation-only checks
  without access to the original execution machine.
\end{enumerate}

\section*{3. Pre-run blockers to close}

These items should be fixed before spending time on the full prospective run.
Otherwise a successful computation can still produce evidence that cannot be
used safely in the paper.

\subsection*{3.1 Correct the QUBO-to-Ising and QAOA objective convention}

\blocked.  The current QAOA primer writes the QUBO coefficients directly as
Pauli coefficients and says the expectation is maximized.  The paper must
instead derive the mapping used by the implementation.  For a minimization
QUBO
\[
 f_Q(z)=c_0+\sum_i q_i z_i+\sum_{i<j}q_{ij}z_i z_j,
 \qquad z_i=\frac{1-Z_i}{2},
\]
the manuscript should display
\[
 H_Q=\kappa I+\sum_i h_i Z_i+\sum_{i<j}J_{ij}Z_iZ_j,
 \quad J_{ij}=\frac{q_{ij}}{4},
 \quad h_i=-\frac{q_i}{2}-\frac14\sum_{j\ne i}q_{ij},
\]
with the corresponding constant $\kappa$.  It must then state unambiguously
whether the classical optimizer minimizes $\langle H_Q\rangle$ or maximizes
$\langle -H_Q\rangle$.  The bit-order convention must be verified against
Qiskit's returned samples with a small exhaustive unit test.

Required changes:
\begin{itemize}
  \item replace the QAOA Hamiltonian paragraph in
  Section ``Quantum solvers'' of \file{main_2.tex};
  \item add a unit test comparing the classical QUBO value and diagonal
  Hamiltonian energy for every bitstring of a small instance, including the
  additive constant;
  \item archive the objective-sense and endianness conventions in each
  prospective artifact.
\end{itemize}

\subsection*{3.2 Make the QRAO description implementation-exact}

\blocked.  The final text must describe the actual variable-to-Pauli
assignment, interaction-graph coloring or grouping rule, encoded Hamiltonian,
ansatz, optimizer, and \texttt{MagicRounding} output produced by the installed
QRAO version.  Compression must be reported from the archived mapping, not
explained by an untested causal story.  Add a machine-readable field for every
variable's qubit and Pauli assignment, and retain the rounded samples used for
each score.  If the installed library cannot export these data reliably, keep
the QRAO discussion descriptive and report only the observed qubit count and
objective distribution.

\subsection*{3.3 Score raw and projected topologies separately}

\blocked.  A cycle-safe projection is an engineering post-processing step; it
must not silently replace a solver's raw bitstring.  Before the run, extend the
prospective artifact schema so each trial contains:
\begin{itemize}
  \item the raw sampled bitstring and raw QUBO objective;
  \item raw connected/radial/component-count/cycle-rank flags;
  \item the projected bitstring, ordered repair actions, and whether any
  action was required;
  \item the projected bitstring's QUBO objective and objective change;
  \item fixed-topology SOC and nonlinear AC validation for the projected
  topology; and, where computationally practical, for the raw topology;
  \item an explicit distinction among ``solver output,'' ``projected
  operational candidate,'' and ``exact hard-radial reference.''
\end{itemize}

Primary QUBO-quality comparisons should use raw outputs.  Engineering
feasibility comparisons may use projected candidates, but must report the
projection rate and cost of repair.

\subsection*{3.4 Complete generated numerical prose and validation tables}

\blocked.  The prospective pipeline currently generates two tables and one
figure, but the abstract and conclusion also need generated macros.  Add:
\begin{itemize}
  \item \file{koshi_admm_qaoa/generated/post_contingency_numbers.tex},
  containing only validated headline macros;
  \item \file{koshi_admm_qaoa/generated/post_contingency_validation_table.tex},
  containing per-method topology, SOC status, maximum cone slack,
  angle-recovery residual, maximum voltage/thermal violation, nonlinear
  convergence, and validated-loss summaries;
  \item consistency checks that fail if \file{main_2.tex} contains a stale
  hard-coded headline value or references a missing generated file;
  \item hashes for both files in
  \file{koshi_admm_qaoa/results/post_contingency_manifest.json}.
\end{itemize}

\subsection*{3.5 Decide the hardware scope before submission}

Choose one of two mutually exclusive paths.

\begin{itemize}
  \item \textbf{Recommended: hardware excluded.} Retain the definitions and evidence
  contract as reproducibility guidance, state that no validated hardware
  package is reported, and make no performance, noise, mitigation, or
  feasibility claim from the old jobs.
  \item \textbf{Hardware included.} Implement a campaign-level aggregator before
  submitting jobs.  It must combine at least the predeclared number of matched
  baseline/noise-management pairs across at least three calibration dates,
  validate every package, preserve per-job results, and calculate paired
  uncertainty without treating shots from one job as independent device
  replications.
\end{itemize}

\section*{4. Required external regeneration}

\external.  A full external run is strictly required before the manuscript
can report post-contingency results from continuous-model-v2.  The present
workspace does not contain the required complete Qiskit/CVXPY execution
environment or the resulting prospective artifacts.

Run the following from the repository root in the documented, frozen
environment after the pre-run blockers above are resolved:

\begin{verbatim}
python koshi_admm_qaoa/study_protocol.py
python koshi_admm_qaoa/artifact_pipeline.py --check

python koshi_admm_qaoa/benchmark.py --post-contingency
python koshi_admm_qaoa/run_post_contingency_sensitivity.py
python koshi_admm_qaoa/make_post_contingency_admm.py

python koshi_admm_qaoa/post_contingency_pipeline.py
python koshi_admm_qaoa/post_contingency_pipeline.py --check

python -m unittest discover -s koshi_admm_qaoa/tests -p 'test_*.py'
\end{verbatim}

The required source artifacts are:
\begin{itemize}
  \item \file{koshi_admm_qaoa/results/post_contingency_v1.json};
  \item \file{koshi_admm_qaoa/results/post_contingency_sensitivity_v1.json};
  \item \file{koshi_admm_qaoa/results/admm_post_contingency_v1.json}.
\end{itemize}

The required derived products are:
\begin{itemize}
  \item \file{koshi_admm_qaoa/generated/post_contingency_results_table.tex};
  \item \file{koshi_admm_qaoa/generated/post_contingency_admm_table.tex};
  \item \file{koshi_admm_qaoa/generated/post_contingency_validation_table.tex};
  \item \file{koshi_admm_qaoa/generated/post_contingency_numbers.tex};
  \item \file{koshi_admm_qaoa/figures/post_contingency_objective_gaps.png};
  \item \file{koshi_admm_qaoa/results/post_contingency_manifest.json}.
\end{itemize}

The run is unacceptable if any trial count is short, any protocol hash
differs, any expected seed is absent or duplicated, any validation field is
missing, a JSON file contains non-finite values, or the post-contingency
manifest check fails.

\subsection*{Optional hardware execution}

Only if the hardware-included path is chosen, first run the local self-test,
then execute the repeated QPU campaign.  A single example package is created
by:

\begin{verbatim}
python koshi_admm_qaoa/phase3_hardware.py --selftest --shots 4096
python koshi_admm_qaoa/phase3_hardware.py --shots 4096

python koshi_admm_qaoa/hardware_evidence.py \
  koshi_admm_qaoa/results/hardware_YYYYMMDDTHHMMSSZ \
  --write-tex koshi_admm_qaoa/generated/hardware_evidence_tables.tex
\end{verbatim}

The timestamped directory name must be replaced by the actual output.  This
validates one evidence package; it does not by itself satisfy the repeated
campaign requirement.  Do not insert the generated hardware tables until the
campaign-level validator also passes.

\section*{5. Artifact-to-manuscript integration map}

\begin{longtable}{p{0.31\textwidth}p{0.27\textwidth}p{0.34\textwidth}}
\toprule
\textbf{Artifact} & \textbf{Manuscript destination} & \textbf{Permitted use} \\
\midrule
\endfirsthead
\toprule
\textbf{Artifact} & \textbf{Manuscript destination} & \textbf{Permitted use} \\
\midrule
\endhead
\file{generated/post_contingency_numbers.tex} & Preamble, immediately after
\file{generated/paper_numbers.tex} & Headline values in the abstract,
Results, Discussion, and Conclusion.  No hand-entered duplicates. \\
\addlinespace
\file{generated/post_contingency_results_table.tex} & Results, immediately
after ``Prospective post-contingency result status'' & Raw-gap distributions,
exact-QUBO hit fractions, and clearly labelled topology/validation rates. \\
\addlinespace
\file{generated/post_contingency_admm_table.tex} & New Results subsection
``Post-contingency ADMM convergence'' & Both residuals, termination status,
final raw/projected topology, projection actions, and current-model
validation. \\
\addlinespace
\file{generated/post_contingency_validation_table.tex} & New Results
subsection ``Physical validation of projected candidates'' & SOC status,
cone tightness, angle recovery, thermal/voltage residuals, nonlinear
power-flow convergence, and loss. \\
\addlinespace
\file{figures/post_contingency_objective_gaps.png} & Results, next to the
primary prospective table & Per-seed or distributional raw objective gaps;
the caption must state the scenario, trial count, and interval definition. \\
\addlinespace
\file{generated/hardware_evidence_tables.tex} & End of ``Hardware evidence
status,'' only after campaign validation & Same-instance expectation, mode,
best sampled energy, exact optima, hit fraction, uncertainty, circuit
metadata, and candidate validation. \\
\addlinespace
Legacy generated files and \file{results/benchmark.json} & Audit appendix or
secondary scaling subsection & Historical base-case surrogate-QUBO audit
only; never labelled post-contingency or current-model evidence. \\
\bottomrule
\end{longtable}

After the new macro file exists, the preamble of \file{main_2.tex} should
contain:

\begin{verbatim}
% Generated by artifact_pipeline.py; do not edit by hand.
\newcommand{\BenchmarkSizeSet}{\{4,6,8,10,12\}}
\newcommand{\MinQuboCouplings}{6}
\newcommand{\MaxQuboCouplings}{66}
\newcommand{\QAOApOneTwelveRatio}{0.987}
\newcommand{\QAOApOneTwelveHitFraction}{0.33}
\newcommand{\ExactQuboMaxTimeMs}{44}
\newcommand{\QAOApOneMinTime}{0.7}
\newcommand{\QAOApOneMaxTime}{1.7}
\newcommand{\NoisyQaoaMaxTime}{96}
\newcommand{\QraoReportedCompression}{1.00\times}
\newcommand{\RadialInfeasibleSizeSet}{\{4,6,8\}}
\newcommand{\PrefixBaselineMaxTimeMs}{0.200}

% Generated by post_contingency_pipeline.py; do not edit by hand.
\newcommand{\PostContingencyRunCount}{30}
\newcommand{\PostContingencyBestGapMethod}{QAOA p1 noiseless}
\newcommand{\PostContingencyBestMedianGap}{0.500}
\newcommand{\PostContingencyBestMedianGapCILower}{0.250}
\newcommand{\PostContingencyBestMedianGapCIUpper}{0.682}
\newcommand{\PostContingencyBestHitMethod}{QAOA p2 noiseless}
\newcommand{\PostContingencyBestHitFraction}{0.467}
\newcommand{\PostContingencyBestGapRawConnectedRate}{0.000}
\newcommand{\PostContingencyBestGapRawConnectedPercent}{0.0}
\newcommand{\PostContingencyBestGapRawRadialRate}{0.000}
\newcommand{\PostContingencyBestGapRawRadialPercent}{0.0}
\newcommand{\PostContingencyBestGapProjectionRate}{1.000}
\newcommand{\PostContingencyBestGapProjectionPercent}{100.0}
\newcommand{\PostContingencyBestGapSocRate}{0.000}
\newcommand{\PostContingencyBestGapSocPercent}{0.0}
\newcommand{\PostContingencyBestGapAcRate}{1.000}
\newcommand{\PostContingencyBestGapAcPercent}{100.0}
\newcommand{\PostContingencyExactOptimumCount}{1}
\newcommand{\PostContingencySensitivityVariantCount}{11}
\newcommand{\PostContingencySensitivityMinMedianGap}{0.000}
\newcommand{\PostContingencySensitivityMaxMedianGap}{2.167}

\end{verbatim}

After the prospective-status paragraph in the Results section, insert the
generated results table, the validation table, and a figure environment that
loads \file{koshi_admm_qaoa/figures/post_contingency_objective_gaps.png}.
Insert the ADMM table in its own immediately following subsection.  The
pipeline should own captions containing numerical facts whenever feasible.

\section*{6. Section-by-section manuscript revision}

\subsection*{Title and abstract}

Restore ``post-contingency'' to the title only if the complete prospective
artifact passes validation and post-contingency results become the paper's
primary evidence.  The abstract should then state:
\begin{itemize}
  \item that the branch-3 Basantapur--Inaruwa circuit-A outage is a
  predeclared \emph{hypothetical} N$-1$ scenario, not a historical event;
  \item the full 13-switch decision set and exact branch ordering;
  \item one or two generated, uncertainty-qualified primary findings;
  \item the separate rates for raw radiality/connectivity, projection, SOC
  validation, and nonlinear AC validation;
  \item that no quantum advantage is observed or claimed.
\end{itemize}
Remove all stale legacy loss values and avoid ``AC-true,'' ``ground truth,''
and unqualified ``success probability.''

\subsection*{Introduction and contribution list}

Replace broad novelty claims with verifiable, bounded language.  The
contributions should distinguish: source-derived topology; estimated
operating point; heuristic historical QUBO; exact topology audit; corrected
continuous model; and auditable prospective protocol.  A result should not be
listed as a contribution until its artifact exists.

\subsection*{System data and continuous formulation}

Retain the corrected complete equations and add a provenance table that labels
each network quantity as source-reported, derived, or assumed.  In particular,
document transformer tap assumptions, branch ratings, conductor assumptions,
load and generation estimates, the 100-MVA base, voltage bases, objective
units, all Big-$M$ constants, and solver tolerances.  State that the nonlinear
check is a fixed-injection power flow, not a nonlinear AC-OPF and not proof of
global AC optimality.

\subsection*{QUBO, radiality, QAOA, and QRAO}

Complete the Ising/objective-sense correction in Section 3.1 of this plan.
Keep the historical cycle and anti-islanding expressions labelled as
heuristic pairwise surrogates and retain the exact flow formulation as the
hard radial reference.  Report exhaustive false-positive and false-negative
behavior beside any surrogate result.  Update the QRAO primer only from the
archived mapping and installed implementation; do not infer why compression
did or did not occur from density alone.

\subsection*{Benchmark methodology}

Preserve the protocol table and exact switch-order table.  Add the execution
environment fingerprint: Python and package versions, solver versions,
operating system, CPU, available memory, and whether times are wall, CPU, or
QPU-access times.  State the contingency, fixed-open/fixed-closed sets, all
penalties, $\rho$ schedule, initialization, stopping criteria, shots,
optimizer settings, seed allocation, and projection rule in one reproducible
location.

Use the offset-invariant gap $f_Q(z)-f_Q^*$ as the primary QUBO metric.
Approximation ratios may be supplementary only if their additive-offset
dependence is explained.  State in advance how ties and degenerate exact
optima are handled.

\subsection*{Results}

Lead with the prospective post-contingency experiment.  Present raw solver
quality first, topology feasibility second, projection effects third, and
continuous validation fourth.  Do not average together failed and successful
physical validations by assigning arbitrary infinite losses.  Report
validation denominators explicitly.

Move the retained base-case size sweep to an audit subsection or appendix so
readers cannot mistake it for the primary post-contingency evidence.  Keep
method rankings scoped to the measured quantity: QUBO objective, runtime,
topology feasibility, or validated physical loss are different comparisons.

\subsection*{Hardware section}

If hardware remains excluded, keep only a short evidence-status paragraph and
move the detailed contract to an appendix or repository documentation.  If
hardware is included, report expectation, standard error, modal sample,
best-energy sampled bitstring, all exact optimum bitstrings, exact-optimum hit
fraction, circuit depth/two-qubit count, backend/calibration metadata, and
same-size validation.  Make only descriptive statements unless the repeated
paired campaign supports an uncertainty-qualified comparison.

\subsection*{Discussion and conclusion}

Rewrite both sections after the generated results are final.  Separate what
the experiment observes from what the formulation guarantees.  Discuss the
cost and frequency of topology projection, SOC tightness failures, nonlinear
recovery failures, estimated data, assumed unit transformer taps, small
problem size, simulator/hardware limits, and sensitivity to penalties.

The conclusion should contain no number absent from a generated macro and no
causal statement absent from a controlled comparison.  It should not claim
AC optimality, hard radiality from the heuristic QUBO, post-contingency
validation from the legacy artifact, or quantum advantage.

\subsection*{Appendices and code/data availability}

Retain the legacy fixed-topology table with a prominent
continuous-model-v1/stale label.  Add appendices for the complete continuous
model, variable/branch ordering, exact radiality formulation, residual
definitions, and evidence schemas if page limits permit; otherwise place them
in a clearly versioned supplement.

The availability statement must name the authoritative files, the archived
commit, license, environment specification, and permanent release identifier.
Archive the exact submission revision with a DOI or equivalent immutable
record rather than citing a moving default branch alone.

\section*{7. Statistical and reporting acceptance rules}

\begin{itemize}
  \item Preserve all 30 predeclared, seed-matched trials per stochastic method.
  Any failure remains in the denominator and is classified by failure mode.
  \item Report median, interquartile range, and the predeclared paired-seed
  bootstrap 95\% interval for continuous gaps; retain the bootstrap seed and
  10,000-resample rule from the protocol.
  \item Report Wilson 95\% intervals for binary rates, including exact-QUBO
  hits, connected/radial outputs, projection use, SOC validation, and
  nonlinear AC validation.
  \item Do not treat shots as independent optimizer runs or independent QPU
  calibrations.  State the unit of replication for every uncertainty estimate.
  \item Treat penalty and $\rho$ sweeps as one-factor-at-a-time sensitivity
  analyses, not causal ablations.  Report all configured levels, including
  failed or unfavorable runs.
  \item Report runtime hardware, software, inclusion boundaries, and whether
  compilation, transpilation, queueing, and validation time are included.
  \item Report exact-QUBO objective ties as sets of optima.  Do not compare a
  modal sample to one arbitrarily selected optimum if the optimum is
  degenerate.
\end{itemize}

\section*{8. External verification still required}

These tasks require source checking or external records rather than another
local code run.

\begin{enumerate}
  \item \textbf{Network provenance.}  Verify every branch endpoint, circuit,
  length, voltage class, conductor/impedance assumption, branch rating,
  transformer parameter, load, and generation value against the cited NEA and
  grid-company sources.  Publish a source/derived/estimated provenance table.
  \item \textbf{Operating scenario.}  Explain why the hypothetical branch-3
  outage is a meaningful stress case.  Report the intact and post-outage base
  states under the corrected model.  Do not imply that the outage occurred
  historically unless an operational record establishes that fact.
  \item \textbf{Novelty.}  Perform a current, systematic literature search for
  transmission reconfiguration, radiality QUBOs, ADMM--QAOA/QRAO power-system
  work, and nonlinear validation.  Replace ``none'' and ``nearly all'' claims
  with bounded statements supported by the search protocol.
  \item \textbf{Bibliography.}  Verify every provisional metadata field and
  remove internal notes.  In particular, resolve the note attached to
  \texttt{morstyn2023annealing} and the author-name note attached to
  \texttt{bhandari2020inps}; remove duplicate or project-internal bibliography
  records if they are not cited.
  \item \textbf{Hardware, if included.}  Verify job IDs against exported job
  results, backend target/properties snapshots, calibration timestamps,
  transpiled circuits, primitive options, and raw counts.  Repository strings
  alone are not evidence of an executed job.
\end{enumerate}

\section*{9. Final submission checks}

Perform the following only after the scientific artifacts and prose are
frozen.

\begin{enumerate}
  \item Run all unit tests, the legacy artifact check, the prospective artifact
  check, and any hardware campaign check.  Save the check outputs with the
  release metadata.
  \item Search the manuscript and supplement for
  \texttt{TODO}, \texttt{PENDING}, provisional notes, ``AC-true,'' stale
  legacy losses, old switch counts, old sample-probability language, and
  unsupported hardware wording.
  \item Compile \file{main_2.tex} from a clean auxiliary-file state.  Treat
  undefined citations/references, missing generated inputs, multiply defined
  labels, and serious overfull boxes as failures.
  \item Visually inspect every page, especially two-column tables, figure
  legends, equation breaks, appendix inputs, URLs, and bibliography
  capitalization.
  \item Confirm that every figure caption names the experiment and evidence
  scope accurately.  Do not claim a panel that is absent from the file.
  \item Add the journal-required declarations: author contributions, funding,
  conflicts of interest, data/code availability, acknowledgments, and any
  generative-AI disclosure required by the target journal.
  \item Create an immutable release containing source, strict raw artifacts,
  generated products, manifests, environment file, tests, README execution
  instructions, commit hash, and DOI.  Verify that the paper cites that exact
  release.
\end{enumerate}

\section*{10. Submission gate}

The manuscript is ready for submission only when every statement below is
true.

\begin{itemize}
  \item[$\square$] The QUBO-to-Ising mapping, objective sense, bit order, and
  QRAO mapping match executable unit tests and archived metadata.
  \item[$\square$] Raw and projected outputs are separately archived and
  reported.
  \item[$\square$] The prospective primary, sensitivity, and ADMM artifacts
  pass strict validation against the frozen protocol hash.
  \item[$\square$] Current-model SOC diagnostics and nonlinear AC recovery are
  reported with denominators and residuals; no stale loss is presented as a
  new result.
  \item[$\square$] Every headline number is generated from the validated
  prospective artifact and covered by the manifest.
  \item[$\square$] The manuscript clearly separates heuristic QUBO behavior,
  exact radiality, topology projection, SOC relaxation, and nonlinear
  power-flow validation.
  \item[$\square$] Hardware is either excluded without residual quantitative
  claims or supported by a validated repeated campaign.
  \item[$\square$] Data provenance, novelty, and bibliographic metadata have
  been externally verified.
  \item[$\square$] A clean compile and visual inspection pass, and the exact
  submission revision is archived immutably.
\end{itemize}

Until these gates pass, the scientifically safe manuscript is the current
audited version: it may describe the corrected methods and the limitations of
the legacy record, but it should not restore quantitative post-contingency,
current-model AC, ADMM-convergence, or hardware-performance claims.

\end{document}